% Work/MAT | 1. rocnik | 23. hodina
% Presne hodnoty sin/cos z geometrie; tg/cotg ze vztahu sin/cos, cos/sin.
% Zdroj: ucebnice Goniometricke funkce.tex, sekce Tabulkove hodnoty;
% Realisticky.cz 040315 a 010121; ilustrace a ulohy vlastni.
\documentclass[11pt,a4paper]{article}
% Spolecny styl pro ucitelske pripravy MAT (Work/MAT)
% Verze 1.1, 2026-10-23
\NeedsTeXFormat{LaTeX2e}
\ProvidesPackage{priprava}[2026/10/23 v1.1 Sablona ucitelskych priprav MAT]

% Zakladni typografie a cestina
\RequirePackage[T1]{fontenc}
\RequirePackage[utf8]{inputenc}
\RequirePackage[czech]{babel}
\RequirePackage{lmodern}
\RequirePackage{microtype}
\RequirePackage{amsmath,amssymb}
\RequirePackage{xcolor}
\RequirePackage{enumitem}
\RequirePackage[most]{tcolorbox}
\RequirePackage{tabularx,booktabs,array}
\RequirePackage{fancyhdr}
\RequirePackage{needspace}
\RequirePackage{geometry}
\RequirePackage[hidelinks]{hyperref}

\geometry{a4paper,top=17mm,bottom=18mm,left=23mm,right=23mm}
\definecolor{notegray}{gray}{0.39}
\setlength{\parindent}{0pt}
\setlength{\parskip}{2.5pt plus 1pt minus 0.5pt}
\setlist[itemize]{topsep=2pt,itemsep=1pt,parsep=0pt,leftmargin=1.45em}
\setlist[enumerate]{topsep=2pt,itemsep=2pt,parsep=0pt,leftmargin=1.65em}
\setcounter{secnumdepth}{0}
\renewcommand{\arraystretch}{1.12}

% Zahlavi: vlevo rocnik | globalni tema, vpravo cislo: nazev hodiny
% Zapati zustava zcela prazdne.
\pagestyle{fancy}
\fancyhf{}
\renewcommand{\headrulewidth}{0.35pt}
\renewcommand{\footrulewidth}{0pt}
\setlength{\headheight}{14pt}
\newcommand{\nastavPripravu}[4]{%
  \fancyhead[L]{\small #1\enspace\textbar\enspace #2}%
  \fancyhead[R]{\small #3: #4}%
}

% Minimalisticke cernobile boxy. Kazda definice zvlast.
\tcbset{commonlesson/.style={enhanced,breakable=false,colback=white,colframe=black,
  colbacktitle=white,coltitle=black,boxrule=0.45pt,
  arc=0.8mm,outer arc=0.8mm,left=3mm,right=3mm,top=1.6mm,bottom=1.6mm,
  toptitle=1mm,bottomtitle=0.7mm,fonttitle=\bfseries,
  attach boxed title to top left={xshift=2.5mm,yshift=-1.15mm},
  boxed title style={colback=white,colframe=white,boxrule=0pt,left=1mm,right=1mm}}}
\newtcolorbox{zadaniBox}[1][]{commonlesson,title={Zadání},#1}
\newtcolorbox{prikladBox}[1][]{commonlesson,title={Příklad},#1}
\newtcolorbox{definiceBox}[1][]{commonlesson,title={Definice},#1}

% Spolecny box pro zadani a jeho reseni, bez mezery mezi castmi.
% Pouziti:
%   \begin{prikladReseniBox}
%     Zadani ulohy.
%     \tcblower
%     \begin{align*} ... \end{align*}
%   \end{prikladReseniBox}
% Spodni cast ma jemny sedy podklad a delici linku; lze v ni sazet
% vice radku matematiky. Obe casti maji jen jeden spolecny ramecek.
\tcbset{solutionlesson/.style={
  commonlesson,
  enhanced jigsaw,
  breakable,
  colbacklower=black!3,
  segmentation style={solid,draw=black!25,line width=0.35pt},
  fontlower=\small\color{notegray},
  before lower={\noindent\textbf{Řešení:}\enspace},
  leftlower=3mm,rightlower=3mm
}}
\newtcolorbox{prikladReseniBox}[1][]{solutionlesson,title={Příklad},#1}
\newtcolorbox{zadaniReseniBox}[1][]{solutionlesson,title={Zadání},#1}

% Pomucky pro autora pripravy: sedive poznamky; \reseni pro starsi pripravy
\newcommand{\poznamka}[1]{\par\begingroup\small\color{notegray}\itshape
  \textnormal{\textbf{Poznámka pro učitele.}}\enspace #1\par\endgroup}
\newcommand{\reseni}[1]{\par\begingroup\small\color{notegray}
  \textbf{Očekávané řešení:}\enspace #1\par\endgroup}
\newcommand{\blok}[1]{\needspace{5\baselineskip}\subsection*{#1}}
\newcommand{\vyrok}[1]{\uv{#1}}

\usepackage{tikz}
\usetikzlibrary{calc}
\nastavPripravu{1. ročník}{Goniometrie}{23}{Přesné hodnoty a kalkulačka}
\begin{document}
\blok{Dnes hodnoty odvodíme, ne jen zapamatujeme (0.--3. minuta)}
V minulé hodině jsme zavedli čtyři poměry pro ostrý úhel. Dnes zjistíme,
\emph{proč} vycházejí známé hodnoty pro $30^\circ$, $45^\circ$, $60^\circ$.
Využijeme dva trojúhelníky, které můžeme přesně sestrojit.

\blok{Z poloviny čtverce získáme $45^\circ$ (3.--12. minuta)}
\begin{prikladReseniBox}[title={P1: rovnoramenný pravoúhlý trojúhelník}]
Dvě odvěsny mají délku $1$. Vypočti přeponu a z toho
$\sin45^\circ$ a $\cos45^\circ$.
\tcblower
Přepona je úhlopříčkou čtverce o straně $1$. Oba ostré úhly
jsou stejné, proto každý měří $45^\circ$.
\begin{center}
\begin{tikzpicture}[scale=.70, every node/.style={font=\small}]
\coordinate (A) at (0,0); \coordinate (C) at (4,0);
\coordinate (B) at (4,4); \coordinate (D) at (0,4);
\draw[densely dashed,draw=black!55] (A)--(D)--(B);
\draw[very thick] (A)--node[below]{$1$}(C)--node[right]{$1$}(B)
  --node[above,sloped]{$\sqrt2$}(A);
\draw (3.65,0)--(3.65,.35)--(4,.35);
\draw (.8,0) arc[start angle=0,end angle=45,radius=.8];
\draw (4,3.2) arc[start angle=-90,end angle=-135,radius=.8];
\node at (1.22,.33){$45^\circ$}; \node at (2.97,3.05){$45^\circ$};
\node[below left] at (A){$A$}; \node[below right] at (C){$C$};
\node[above right] at (B){$B$};
\end{tikzpicture}
\end{center}
$1^2+1^2=c^2$, takže $c=\sqrt2$ a
\[
\boxed{\sin45^\circ=\cos45^\circ=\frac1{\sqrt2}=\frac{\sqrt2}{2}}.
\]
\end{prikladReseniBox}

\blok{Výška rovnostranného trojúhelníku: $30^\circ$ a $60^\circ$ (12.--23. minuta)}
\begin{prikladReseniBox}[title={P2: rozděl trojúhelník výškou}]
Rovnostranný trojúhelník má stranu délky $2$. Sestroj výšku.
Odvoď délky odvěsen a hodnoty $\sin$ a $\cos$ pro $30^\circ$ a $60^\circ$.
\tcblower
Výška je zároveň těžnicí a půlí vrcholový úhel.
Dostaneme dva trojúhelníky $30^\circ$--$60^\circ$--$90^\circ$.
\begin{center}
\begin{tikzpicture}[scale=.91,every node/.style={font=\small}]
  \coordinate (A) at (0,0);\coordinate (B) at (3.6,0);\coordinate (C) at (1.8,3.1177);
  \draw[very thick] (A)--node[below]{$2$}(B)
    --node[right,sloped]{$2$}(C)--node[left,sloped]{$2$}(A);
  \node[below left] at (A){$A$}; \node[below right] at (B){$B$};
  \node[above] at (C){$C$};
  \node at (1.8,-.98){\emph{rovnostranný}};
  \node at (4.55,1.45){$\Longrightarrow$};
  \begin{scope}[xshift=5.5cm]
  \coordinate (P) at (0,0);\coordinate (Q) at (3.6,0);
  \coordinate (R) at (1.8,3.1177);\coordinate (S) at (1.8,0);
  \fill[black!9] (P)--(S)--(R)--cycle;
  \draw[very thick] (P)--(Q)--(R)--cycle;
  \draw[dashed,thick] (R)--(S);
  \draw (1.53,0)--(1.53,.27)--(1.8,.27);
  \node[below] at (.82,-.08){$1$};\node[below] at (2.75,-.08){$1$};
  \node[left] at (.8,1.50){$2$};
  \node[right] at (2.15,1.75){$v=\sqrt3$};
  \draw (.65,0) arc[start angle=0,end angle=60,radius=.65];
  \node at (.95,.40){$60^\circ$};
  \draw (1.8,2.5) arc[start angle=-90,end angle=-120,radius=.62];
  \node at (1.32,2.38){$30^\circ$};
  \node[below left] at (P){$A$};\node[below right] at (Q){$B$};
  \node[above] at (R){$C$};\node[above left] at (S){$D$};
  \node at (1.8,-.98){\emph{rozdělení výškou}};
  \end{scope}
\end{tikzpicture}
\end{center}
Z Pythagorovy věty: $v^2+1^2=2^2$, tedy $v=\sqrt3$.
Proti $30^\circ$ leží odvěsna $1$, proti $60^\circ$ výška $\sqrt3$:
\[
\boxed{\sin30^\circ=\tfrac12,\quad\cos30^\circ=\tfrac{\sqrt3}{2}},
\qquad
\boxed{\sin60^\circ=\tfrac{\sqrt3}{2},\quad\cos60^\circ=\tfrac12}.
\]
\end{prikladReseniBox}

\newpage
\blok{Hodnoty $0^\circ$, $90^\circ$ a tangens s kotangensem (23.--30. minuta)}
\textbf{Krajní úhly:} v opravdovém pravoúhlém trojúhelníku nemůže být
jiný ostrý úhel roven $0^\circ$ nebo $90^\circ$. Pro tyto dvě hodnoty
použijeme stručně \emph{jednotkovou kružnici}: bod pro $0^\circ$ má
souřadnice $(1,0)$ a pro $90^\circ$ souřadnice $(0,1)$.
Kosinus je vodorovná a sinus svislá souřadnice.
\begin{center}
\begin{tikzpicture}[scale=.74]
\draw[->] (-2,0)--(2,0) node[right]{$x$};
\draw[->] (0,-1.9)--(0,2) node[above]{$y$};
\draw[thick] (0,0) circle(1.45);
\fill (1.45,0) circle(2pt) node[below right]{$0^\circ\ (1;0)$};
\fill (0,1.45) circle(2pt) node[above right]{$90^\circ\ (0;1)$};
\node[below left] at (0,0){$O$};
\end{tikzpicture}
\end{center}
Z definic z 22. hodiny pro \emph{ostré} úhly a z tohoto geometrického
rozšíření platí (vždy, když jmenovatel není nula):
\[
\tan\alpha=\frac{\sin\alpha}{\cos\alpha},\qquad
\cot\alpha=\frac{\cos\alpha}{\sin\alpha}.
\]
\textbf{Odvození pro učitele (na tabuli stačí jeden sloupec):}
\[
\begin{aligned}
\tan30^\circ&=\frac{1/2}{\sqrt3/2}=\frac{\sqrt3}{3}, &
\cot30^\circ&=\frac{\sqrt3/2}{1/2}=\sqrt3,\\
\tan45^\circ&=\frac{\sqrt2/2}{\sqrt2/2}=1, &
\cot45^\circ&=1,\\
\tan60^\circ&=\frac{\sqrt3/2}{1/2}=\sqrt3, &
\cot60^\circ&=\frac{1/2}{\sqrt3/2}=\frac{\sqrt3}{3}.
\end{aligned}
\]
Na krajích: $\tan0^\circ=0/1=0$, $\cot0^\circ$ není definován;
$\tan90^\circ$ není definován, $\cot90^\circ=0/1=0$.
\begin{center}
\renewcommand{\arraystretch}{1.35}
\begin{tabular}{@{}lccccc@{}}\toprule
$\alpha$ & $0^\circ$ & $30^\circ$ & $45^\circ$ & $60^\circ$ & $90^\circ$ \\ \midrule
$\sin\alpha$ & $0$ & $\frac12$ & $\frac{\sqrt2}{2}$ & $\frac{\sqrt3}{2}$ & $1$ \\
$\cos\alpha$ & $1$ & $\frac{\sqrt3}{2}$ & $\frac{\sqrt2}{2}$ & $\frac12$ & $0$ \\
$\tan\alpha$ & $0$ & $\frac{\sqrt3}{3}$ & $1$ & $\sqrt3$ & nedef. \\
$\cot\alpha$ & nedef. & $\sqrt3$ & $1$ & $\frac{\sqrt3}{3}$ & $0$ \\ \bottomrule
\end{tabular}
\end{center}
\poznamka{Trojúhelník vysvětluje přesné hodnoty pro ostré úhly.
Kraje vysvětluje kružnice; netvrď, že jde o přípustné ostré úhly.
Tabulku nenuť celou memorovat hned. Tangens a kotangens ukaž
podílem již vypočtených hodnot.}

\newpage
\blok{Kalkulačka: správný režim, hodnoty a opačná úloha (30.--41. minuta)}
\textbf{Postup:} přepni kalkulačku na \texttt{DEG} (stupně), nikoli
\texttt{RAD}. Kontrola: $\sin30^\circ=0{,}5$.
Pro výpočet hodnoty stiskni např. \texttt{sin}, zadej $37$ a potvrď;
některé kalkulačky vyžadují nejprve číslo.
Chceme-li \emph{určit ostrý úhel} ze známé hodnoty, použijeme
\texttt{SHIFT} + \texttt{sin} (zobrazí \texttt{sin\ensuremath{^{-1}}}),
popř. \texttt{arcsin}; analogicky \texttt{arccos}, \texttt{arctan}.
Nejde o převrácenou hodnotu čísla!
\begin{prikladReseniBox}[title={P3: hodnoty z kalkulačky},breakable=false]
Spočítej $\sin37^\circ$, $\cos53^\circ$. Proč jsou výsledky stejné?
\tcblower
$\sin37^\circ\doteq0{,}6018$, $\cos53^\circ\doteq0{,}6018$.
Úhly se doplňují do $90^\circ$: platí
$\sin\alpha=\cos(90^\circ-\alpha)$ pro ostrý $\alpha$.
\end{prikladReseniBox}
\begin{prikladReseniBox}[title={P4: opačná úloha s kalkulačkou},breakable=false]
Urči ostrý úhel $\alpha$, pro který je $\cos\alpha=0{,}6$.
\tcblower
$\alpha=\arccos(0{,}6)\doteq53{,}13^\circ$.
Pro kontrolu $\cos53{,}13^\circ\doteq0{,}6$ (se zaokrouhlením).
Úlohu řešíme v oboru ostrých úhlů.
\end{prikladReseniBox}
\blok{Výstup a domácí práce (41.--45. minuta)}
\begin{zadaniReseniBox}[title={S1: krátká kontrola},breakable=false]
\textbf{a)} Bez kalkulačky $\tan45^\circ+\cot60^\circ$.
\quad \textbf{b)} Jak poznáš na kalkulačce omylem zapnuté radiány?
\tcblower
a) $1+\frac{\sqrt3}{3}$. b) Test $\sin30^\circ$ nedá $0{,}5$.
\end{zadaniReseniBox}
\begin{zadaniReseniBox}[title={D1: doma, 10 minut},breakable=false]
\begin{enumerate}[label=\alph*)]
\item Z geometrie odvoď znovu $\sin45^\circ$.
\item Z tabulky vypočti $2\sin30^\circ+\cos60^\circ$.
\item Na kalkulačce vypočti $\tan25^\circ$ a $\arcsin(0{,}6)$.
\end{enumerate}
\tcblower
a) $1/\sqrt2=\sqrt2/2$.
b) $2\cdot\frac12+\frac12=\frac32$.
c) $\tan25^\circ\doteq0{,}4663$, $\arcsin(0{,}6)\doteq36{,}87^\circ$.
\end{zadaniReseniBox}
\poznamka{Rozpis: 3+9+11+7+11+4=45 minut.
Při pomalejší třídě nechej výpočet $\cot$ v tabulce jako hotový výsledek,
ale odvození ponech v přípravě pro vlastní použití.}
\end{document}
